\chapter{Maps Look Objective, but Always Make Choices}
\section{The Map on the Classroom Wall}
First, take up something you see every day but rarely really look at: the world map on your classroom's back wall.

Go closer. At its center are probably the Pacific Ocean and China on its western shore. China sits comfortably just left of center. Europe is squeezed against the left edge, the Americas pushed to the right, and Greenland hangs high above, looking like a white continent almost as large as Africa. You have memorized the seven continents and four oceans before this map, found countries, perhaps located the host of a football match. It hangs quietly like a window, as if this were simply what the world should look like when made smaller.

Now try a small experiment. Imagine passing through Earth's center into a classroom on the other side, perhaps in London or Paris. Its wall also has a world map, but the Atlantic sits in the middle. Europe and Africa occupy the center, the Americas are on the left, Asia is pushed right, and China lies far away near an edge. One Earth, two maps, each centering itself and banishing the other to the margins. Which is wrong?

The answer may be uncomfortable: neither is wrong, and both have arranged things.

This is the chapter's subject. A map seems an utterly honest version of the world: coastlines, borders, mountains, rivers, cities, all clearly marked like a photograph. Yet before any map exists, someone makes a series of choices. Where is the center? Which way is up? What scale? How will a sphere be flattened? What will be shown or omitted? People make these choices for purposes, and the results quietly influence every reader. Maps look objective, but always make choices. Learning to recognize those choices is a basic skill for reading the world.

\section{A Ruler in a Berlin Conference Room}
Now come with me into a scene.

Time: November 1884, winter in Berlin. Place: a conference room in German chancellor Bismarck's residence. The fireplace blazes, cigar smoke hangs in the air, and green baize covers a long table. Representatives of more than a dozen countries sit on either side: Germany, France, Britain, Portugal, Belgium, Spain, and even the distant United States. Notice something astonishing to modern ears. The conference will discuss dividing the entire African continent, yet no African is in the room. Not one African ruler has been invited.

A huge map of Africa hangs on the wall. Its coastline is detailed. For centuries, European merchant ships have traded along the coast and established outposts, repeatedly surveying every bay. Inland, however, great blank spaces remain. European knowledge of the interior's mountains, rivers, deserts, and kingdoms is limited. The blanks seem to say there is little worth drawing.

What happens over the following months changes tens of millions of lives. Delegates gather before the map, argue, bargain, exchange concessions, then take rulers and draw lines across the blanks. One straight line from a river to a mountain, another directly along a parallel of latitude. Nobody asks the people on either side about their ethnic groups, languages, pastures, water sources, or longstanding feuds with neighbors. The ruler does not care. Its advantages are speed, clarity, and the ability to finish a line at the negotiating table.

By the conference's end, Africa's partition is largely settled on paper. Over the next thirty-odd years, those lines reach the ground one by one. Colonial armies enter the interior, plant boundary stones, levy taxes, and build railways. Peoples who never considered themselves one country are enclosed by one border, while communities that have lived together for generations are cut in two. Open a world map today and Africa still bears the conference's handwriting: straight borders following latitude and longitude regardless of terrain. You can almost distinguish them intuitively. Curving borders often follow rivers and mountains and grew from life on the ground; straight ones were usually drawn. At Namibia's northeastern corner, a peculiar narrow arm reaches hundreds of kilometers inland, scarcely shaped like territory at all. It was exchanged on a map by treaty a few years later so that a German colony could reach the Zambezi River.

Remember this room. Half the chapter's question lies hidden in that ruler.

\section{What Is a Map Saying?}
Let us lay out the conflict without rushing to an answer.

One respectable, popular view might be called the photograph view: a map is a reduced photograph of the world. Accuracy defines a good map. More precise coastlines, more consistent scale, smaller errors are always better. Faults reflect inadequate technology. As surveying satellites and global positioning systems fill the sky, maps approach perfect copies. On this account, the Berlin conference's problem was merely inaccurate drawing. Straight lines arose from ignorance; technological progress would naturally make maps more correct.

A sharper view is the choice view: a map never was a photograph and never can be. Earth is spherical and paper flat, so flattening it necessarily sacrifices something. The world is infinitely rich and paper finite, so including something excludes something else. Every map results from choices made by people with purposes, positions, and interests. From this perspective, Berlin's problem was not an inaccurate ruler, but the ruler itself: who was entitled to hold it, who set the criteria, and why those enclosed by its lines could not even attend.

Both sides have facts behind them. The photograph camp can point to satellite imagery: maps today really are orders of magnitude more accurate than two hundred years ago, genuine progress. The choice camp can point to our classroom maps: if the technology and surveying accuracy are the same, why does one center the Pacific and the other the Atlantic? That difference is not error. It is viewpoint.

So the real question is \textbf{what relationship a map has to the world}.

If a map is a photograph, we need better cameras and can leave the rest to technology. If it is a choice, every map deserves questioning: Who made it? For whom? What was sacrificed or enlarged? Who was left outside the paper? These approaches lead to different ways of reading maps and different worlds.

Before reading on, take your own position. Was the Berlin ruler's main problem inaccurate drawing, or denying others their turn to draw?

\section{Worlds Seen from Different Positions}
Place several very different historical world pictures side by side, and something immediately becomes clear: the center and top have always been reserved for what matters most.

Many medieval European world maps put east at the top and Jerusalem at the center. In believers' worlds, the holy city naturally formed the axis. The English word orientation comes from oriens, east, a fingerprint those maps left in language. At roughly the same time, Islamic cartography took another path. In the twelfth century, al-Idrisi made a world map for King Roger II of Sicily, an outstanding map with south at the top. In China's Han tombs at Mawangdui, archaeologists found silk maps more than two thousand years old that also put south above north. For their makers, the ruler sat facing south, so south naturally belonged above. North-up is no heavenly law, merely one of many possibilities that happens to dominate today.

Move now to the Pacific. Before Europeans arrived with paper and compasses, Marshall Islands navigators used maps that astonished Western sailors: frameworks of palm ribs whose curves represented currents and swells, with small shells tied on for islands. These stick charts were not drawn on paper but tied to canoes and held in bodily memory. Navigators felt swells through their skin to judge islands' directions. In 1769, when Captain Cook's ships reached Tahiti, a local priest and navigator named Tupaia joined the voyage. From memory, he described and drew dozens of islands across an oceanic area far wider than European charts then covered. The same Pacific appeared on European maps as large blanks and scattered routes, but on Tupaia's as dense islands and passages. One was an outsider's sea, the other the roads around home.

Now hear different voices. Instead of asking who draws, ask how space arranges human life.

Traders understand this early. On the ancient Silk Roads, caravans could not simply go wherever they pleased. Geography led them. Oases set supply points, passes set routes, and paths formed along both edges of the Taklamakan Desert. Venice and Genoa competed not for land so much as straits and ports. Whoever controlled a sea route could tax what passed through. Space entered trade's accounts directly.

Soldiers understand it too. At Thermopylae in 480 BCE, the Greek alliance held a coastal passage only a few people wide. Persian numerical superiority could not deploy, and the army was delayed three days. China's Hangu and Tong passes repeatedly made their narrow corridors unavoidable in dynastic struggles. In the early twentieth century, the British geographer Halford Mackinder even proposed a spatial explanation of world history: broadly, whoever controlled Eurasia's interior held the throat of the World-Island and could dominate the world. His Heartland theory influenced a century of strategists. You may disagree, but must acknowledge that he read the map as a manual of power.

Migration follows space as well. When sea levels fell during ice ages, a land bridge emerged at the Bering Strait, and people moved from Asia into the Americas. Polynesians followed winds, currents, and stars, island-hopping eastward across generations and making the central Pacific their home. Migration routes are never random lines. They follow water, grasslands, and places where life is possible.

Another voice has remained outside the paper: people placed on maps but never allowed to draw them. When Berlin drew its lines, Africa already had kingdoms, trade routes, pastures, and markets, with its own divisions and memories of land. European maps turned these into blanks. Once the lines reached the ground, the consequences were concrete. Somali people who had moved camels seasonally for generations found themselves divided among five jurisdictions, crossing international borders to reach the same well. A straight border split the Maasai of East Africa. Cattle did not recognize boundary stones, but people could pay in blood for crossing. For them, the map was neither a photograph nor a neutral tool, but paper descending from above to determine where they could go. Those being named had no right to name: cartography's oldest inequality.

Finally, space can become identity. Think of Chinese regional terms: Jiangnan, south of the Yangtze; Saibei, beyond the northern frontier; Guandong, east of the pass; Lingnan, south of the ranges; Zhongyuan, the Central Plains. Each begins as geography but trails imagined character behind it. Jiangnan suggests watery towns and delicacy; Saibei, windblown sand and boldness. The Qinling--Huai River line divides one country into southerners and northerners, who then earnestly dispute whether soft tofu should be sweet or savory. A line drawn long enough on a map enters people's minds and becomes us and them. Berlin's lines grew, decades later, into real countries, passports, and mutual suspicion: strong evidence that maps help make boundaries.

Together, these voices reveal a shared outline. Centers, directions, and borders answer one question: what matters most to the mapmaker? People in different positions never give the same answer.

\section{Thinking Bridge: Four Questions for a Map}
Can we analyze a map with a portable tool rather than intuition alone? Yes. Ask these four questions in order and put any map through the same examination.

\textbf{First: Where are the center and top?}

This asks whose world it is. Pacific-centered or Atlantic-centered? North-up or south-up? Is your city central or marginal? The center and top are a map's most expensive real estate. See who occupies them and you see its position.

\textbf{Second: How was the sphere flattened?}

Earth is a sphere and paper a plane. Flattening a sphere is like peeling and flattening an orange: something must tear or stretch. Mathematicians long ago proved that no completely distance-preserving correspondence exists between curved and flat surfaces. Every flat world map must therefore distort; the question is where and how much. The method of translating sphere into plane is called a projection. The famous Mercator projection, published in 1569 by the Flemish cartographer Gerardus Mercator, preserves angles and shapes at the price of enormous stretching at high latitudes. Greenland looks nearly as large as Africa, although Africa is about fourteen times its area. Other projections followed, including the Peters projection, which caused major debate in the 1970s. It preserves relative areas but stretches and squashes continents into unfamiliar shapes. There is no perfect projection, only different choices of what to sacrifice. Asking which projection was used means asking what the map gives up.

\textbf{Third: What is the scale?}

Scale is the ratio of map distance to real distance, but you can think of it as permitted omission. A 1:10000 map can show your school and the intersection outside your home. On a world map, the whole city becomes a point and your street disappears. Larger scales show more detail and discard less; smaller scales discard more. So ask what is allowed to vanish at this scale.

\textbf{Fourth: What is shown, and what is not?}

This is the most important question. World maps omit house prices, tourist maps factories, navigation maps public safety. Every map leaves things outside its paper, and those silences are often not random.

Use all four questions together. Find four maps of one place, for example China's physical geography, population, languages, and political divisions, and compare them.

The physical map shows three great steps of elevation, with the Tibetan Plateau high above and major rivers flowing east to the sea. It explains why early Chinese civilizations clustered in the Yellow and Yangtze river valleys. The population map shows a line drawn in 1935 by geographer Hu Huanyong, running diagonally from Heihe in Heilongjiang to Tengchong in Yunnan. Southeast of it, about forty percent of the land has consistently held more than ninety percent of the population. More than eighty years later, the line remains broadly effective. The language map exposes the crude simplification of ``Chinese people speak Chinese.'' Tibetan, Uyghur, Mongolian, Zhuang, and many other languages occupy different regions, while differences within Chinese varieties alone exceed those between many European languages. The political map shows provincial boundaries. Some bend with mountains and rivers; others deliberately interlock. Hanzhong Basin, south of the Qinling and geographically closer to Sichuan, belongs to Shaanxi, a deliberate design since the Yuan to prevent a province from relying on natural barriers to become independent.

Four maps, four stories. None lies, and none tells everything. Real map-reading skill is not finding the most accurate one, but identifying the question each answers.

\section{A Three-Thousand-Year-Old Geographical Plan}
Now read a short primary source from one of China's oldest geographical texts: the ``Tribute of Yu'' in the Book of Documents.

Attributed to the flood-tamer Yu the Great, it describes the realm as concentric zones around the royal capital, each five hundred li wide. It describes the innermost zone this way:

\begin{quote}
Five hundred li form the royal domain: within one hundred li, tribute is paid in whole stalks; within two hundred, in cut ears; within three hundred, in grain with stalks and service; within four hundred, in grain; within five hundred, in rice.
\end{quote}
Roughly, the five-hundred-li royal domain pays grain in five grades according to distance. Within one hundred li, deliver stalks and ears together; within two hundred, the ears; within three hundred, coarse grain plus labor; within four hundred, husked rice; within five hundred, polished rice. Nearer the capital, tribute is bulkier and service heavier. Farther away, tribute is lighter and more refined because long-distance transport can carry only the choicest part. Beyond the royal domain, the text lists successive five-hundred-li zones: the lords' domain, the pacification domain, the domain of restraint, and the remote domain. Obligations loosen with distance until the outermost offers only nominal allegiance.

Most scholars date the actual composition to the Warring States period, an ideal plan written in Yu's name rather than a measured survey. That makes it an excellent specimen: a map drawn in words that depicts a wish. Notice its boundaries. Not walls or rivers, but gradients of obligation. Once boundaries are drawn, duties are created. A place assigned to the royal domain ought henceforth to deliver grain. Distant places classified as restrained or remote become margins and wilderness. Maps name boundaries and help make them. People in China were rehearsing this on the page three thousand years ago.

At almost the same time, someone in the Mediterranean doubted every neatly drawn world. In his Histories, Herodotus mocked contemporary maps depicting a perfectly round Earth encircled by water called Ocean. His point was roughly: the land I have seen does not look like that. These mapmakers have never seen it all, yet make the unknown world perfectly orderly. On one side, the Tribute of Yu turns the world into a court; on the other, Herodotus mocks every tidy drawing. Antiquity had already occupied both positions from our opening.

\section{Same Land, Different Maps: Why?}
Return to Berlin's ruler and the two classroom maps. With more pieces in hand, we can offer genuinely different explanations for maps' differences. First separate four matters. What happened, lines drawn in Berlin, high latitudes stretched by Mercator, belongs to evidence. Why it happened belongs to explanation, where views compete. How contemporaries understood it, colonizers believing they were surveying and bringing order, is their self-account. How we evaluate it, whether partition was just, is a value judgment. What follows compares the second level.

\textbf{Explanation one: Geometry imposes limits. The technical account.}

A sphere cannot be flattened. That is mathematics, not conspiracy. Every flat world map must distort, and every small-scale map must omit. Many differences arise from geometrical necessity plus scale limits, regardless of the maker's goodness or badness. This explanation is clean and testable. Why is Greenland so large? Projection: calculate and see. Its limit is equally clear. Geometry demands a sacrifice without choosing which. Mercator sacrifices area to preserve angles; Peters sacrifices shape to preserve area. Geometry does not decide between them. People do. The technical account explains why distortion is necessary, not why this particular distortion was chosen.

\textbf{Explanation two: Mapmakers benefit. The power account.}

Look at who pays for the map. Berlin's map left the interior blank, but blankness was never neutral. In contemporary European opinion, blank meant unowned, awaiting discovery and occupation. A map thus became an invitation to possess. Names work similarly. For whom is the Middle East middle, and for whom east? Only from Europe is it the middle part of the east. This sharp account explains why centers, names, and blanks favor the makers. Its limit is that reducing everything to interest cannot explain popular maps that plainly sacrifice accuracy---

\textbf{Explanation three: Purpose shapes the cut. The functional account.}

---which brings us to this chapter's most important counterexample. In 1931, London Underground employee Harry Beck drew a new network map. Geography was thoroughly sacrificed. Station spacing no longer matched distance; routes followed only horizontal, vertical, and forty-five-degree lines. Central stations spread out and suburban distances compressed. By the photograph camp's standard, the map was wildly wrong. The result? It became the template for subway maps worldwide, copied by countless cities for nearly a century. It honestly answered just one question, how to travel and where to change, and answered it extremely well. Nobody accused Beck of distorting London, because he never pretended to show London's geography.

The counterexample exposes a hasty inference. Not every distortion is deception. A distorted map can be honest if it declares its purpose. The test is not resemblance, but whether it answers the question it claims to answer.

Now state Mercator's case fully. This is worth doing because recent claims that his projection is a colonialist map deliberately belittling tropical countries are unfair to it.

Mercator has three arguments. First, his map has a distinctive mathematical property: a straight line corresponds to a sea route with a constant compass bearing. Sixteenth-century navigators had no satellites, only compasses and the map. Draw the line, read the angle, steer accordingly, and arrive alive. It saved countless sailors over centuries, a marvel in navigation's history. Second, Mercator never intended a portrait of the world. His map's title explicitly states its navigational purpose. Distortion was a functional price, not an ideological plot. Third, electronic nautical charts still use Mercator today because it remains appropriate for navigation.

These arguments stand. But notice the other side of steelmanning. The Mercator map itself is not at fault; the mismatch is. Hang a navigational tool on classroom walls and print it in textbooks as the world's appearance for generations of children, and its distortion becomes their mental world. Northern lands seem large, southern ones small, and large-looking places quietly seem important. The problem is not the map but our choosing the wrong map.

Each account grasps part of the truth. Geometry requires sacrifice, interests influence what is sacrificed, and purpose can make the sacrifice worthwhile. A skilled reader tests all three rather than settling for the first that comes to hand.

\section{Further Challenge: A Map's Silences}
Now for this chapter's weightiest theoretical tool. It comes from the British-born historian of cartography J. B. Harley. In the 1980s, he published essays, most famously ``Deconstructing the Map,'' that turned cartography's perspective around.

Before him, map history was largely a story of progress: ever more precise surveys, fuller drawings, smaller errors, a march toward perfect copying. Harley's point was roughly that this narrative watched what maps showed, not what they omitted, while their real power lay precisely in silence. He saw mapping as two filters. The first is collection. The world is infinitely rich, surveyors can record only some things, and what deserves recording is already a choice. The second is drawing. Which recorded things enter the map, in what type size, in which position? More choices. After both filters, a map appears clean, objective, and neutral, though every part results from judgment.

Reread our examples through this theory and echoes appear everywhere. Berlin's blanks were silence, erasing existing kingdoms, routes, and pastures and turning inhabited into unclaimed. The Tribute of Yu's restrained and remote domains are another silence. The center can name the margins, but the margins cannot name themselves. Beck's subway map is most candid: it silences all of London's surface without ever pretending to depict it. Three silences, three different qualities. This is Harley's desired skill: do not ask only whether a map is objective, but whom it silences and whether it admits doing so.

Push this theory to its limit, then pull it back a little. Every theory has boundaries.

If maps are nothing but power, how do sailors come home alive? A careless chart sends a ship onto rocks; a carefully surveyed one brings it into port. Reefs, not power, decide that difference. Harley attacks the naivety of map as truth, not the idea that some maps are better or more honest than others. The right stance has two levels. Recognize each map's choices and silences, so keep questioning. Recognize that choices differ in quality, so do not lapse into everything-is-conspiracy nihilism. A map can be power's language and an effective model at once. Both statements are needed. The tool is a flashlight, not a verdict: it illuminates silenced corners, but before finding a map guilty, ask what it claims to do, how well it does it, and whom its silence harms.

\section{The Map in Your Pocket and the World Outside}
This three-thousand-year-old question plays out in your pocket every day.

First, navigation maps. Open your phone and a blue dot sits at the center: you. For the first time in cartographic history, everyone has a self-centered world map, a privilege unavailable to Mercator, al-Idrisi, or the Tribute of Yu's author. Yet that centrality is an illusion. You see an interface; algorithms behind it decide what appears. Which businesses come first, which route is recommended, which districts are detailed and which remain gray are questions of commerce and rules, not surveying. Code now guards Harley's second filter. Meanwhile, research suggests that prolonged total dependence on navigation can weaken our own route-finding and route-memory skills. A sense of direction is like a muscle: unused, it shrinks. Marshall Islands navigators read currents with their skin; we read blue dots with our eyes. Who better knows the roads beneath their feet is not obvious.

Next, boundaries. On your city's school-catchment map, one line can make houses on opposite sides of the same street differ sharply in price. The line is not fundamentally different from the Tribute of Yu's rings. Its maker names a boundary; the boundary creates eligibility and interests. Food-delivery areas and shared-bicycle operating zones are electronic domains too: inside is the royal domain, outside the remote domain. Those outside best understand the experience of being named.

Next, statistical maps. On U.S. presidential election nights, television screens show states colored red or blue across vast areas. But land does not vote; people do. Large, sparsely populated states fill the screen, while dense cities nearly vanish. Coloring by land and coloring by population tell different stories about the same election, and politicians know which to choose. Disease maps work similarly: coloring a whole country averages away enormous internal differences. Next time news shows a colored map, ask the four questions, especially the fourth. What unit is colored, and whom does that unit hide?

Finally, identity. Classrooms in every country display maps placing their country centrally or prominently. Geography lessons teach our territory. Maps are the modern state's most ordinary lessons in identity. That is why people care so intensely about every line and every island or reef label. The dispute is never only geography, but the picture of who we are.

\section{Over to You: Which Classroom Map?}
Return to the opening wall.

If your classroom could display only one world map, which would you choose? Pacific-centered, with China prominent? Atlantic-centered, as commonly used in Europe and America? South-up, turning the familiar world upside down? The Peters projection, with true areas but strangely distorted continents? Or a rotating globe instead?

Give an answer and reasons. Before deciding, weigh these arguments again.

You have a case for centering yourself. A map is a perspective anyway, and children everywhere may begin understanding the world from their own position. Starting by finding where I am is nothing shameful.

You have a case for true areas too. A classroom map is not a nautical chart. It teaches relative sizes and positions, so its projection should do least harm to that purpose. We have already seen the cost of mismatched tools with Mercator.

You also know a harder fact: every choice has silences, including the apparently perfect globe. It shows only half of Earth at a time, hiding the other half behind it. Displaying many maps seems comprehensive, but declares each merely one version. For a child seeking certainty, might that be another kind of cruelty?

Then what defines a good map: accuracy, honesty, fairness to the weak, or openly declaring its choices? If maps must choose, can we demand objectivity, or only an account of their choices?

There is no standard answer. But from now on, any map, at school, in the news, on your phone, will make you want to ask those four questions. The map has not changed. You have: instead of merely being led by a map, you can choose a path for reading it.
